% auto-generated by quick_analysis/conviction_summary_tables.py
\begin{table}[H]
\centering
\small
\begin{tabular}{lcccc}
\toprule
 & \multicolumn{2}{c}{High conviction} & \multicolumn{2}{c}{High degree} \\
\cmidrule(lr){2-3}\cmidrule(lr){4-5}
Dataset & Edges & Drop & Edges & Drop \\
\midrule
Amazon Ratings & $32.9$ & $6.43$ & $46.1$ & $5.88$ \\
Citeseer & $37.0$ & $3.27$ & $48.1$ & $1.01$ \\
Cora & $36.3$ & $2.58$ & $46.1$ & $2.68$ \\
Minesweeper & $29.8$ & $17.15$ & $33.2$ & $8.19$ \\
Pubmed & $49.6$ & $1.34$ & $61.5$ & $1.45$ \\
Questions & $52.4$ & $10.31$ & $68.3$ & $12.31$ \\
Roman Empire & $36.5$ & $12.92$ & $44.4$ & $11.48$ \\
Tolokers & $40.4$ & $2.63$ & $65.2$ & $3.17$ \\
\bottomrule
\end{tabular}
\caption{Structural cost against performance cost, per dataset, for pruning at the 30\% budget. ``Edges'' is the percentage of the graph's edges deleted along with the selected nodes; ``Drop'' is the relative drop (\%) in the headline metric, as everywhere else in this section. Selecting by degree deletes the larger fraction of the graph on every dataset, yet on several of them it costs less than selecting by conviction, which is what shows that part of the damage a degree criterion does is structural rather than informational. Conviction score \texttt{last}, values averaged over the ten folds.}
\label{tab:conv_prune_evd_last}
\end{table}
